\section{Conclusion}
\label{sec:conclusion}
This article proposed DA-HELLO, a HELLO scheduler for Q-learning
underwater routers that sends a HELLO on first contact with a new
neighbour, on a change of void status, and, optionally, after a long
silence. Evaluated on two re-implemented hosts, MARV and QELAR, with
state carried either in data headers or only in HELLOs, the results show
that the value of HELLO scheduling depends on how routing state is
disseminated. Where state travels only in HELLOs, DA-HELLO delivers
$2.2$ to $3.0$ percentage points more than a jittered $120\,\mathrm{s}$
periodic schedule on MARV with $43$ to $53\%$ fewer control bytes, and
saves $45$ to $52\%$ of the bytes on QELAR, although non-inferiority in
delivery was not established in its drifting scenarios. Where state
travels in data headers, DA-HELLO matches sending no HELLOs, whereas
periodic and Trickle schedules lose delivery; on QELAR this holds only
with the first-contact reply switched off. The score-change test showed
no demonstrated benefit. These conclusions rest on simulation of two
re-implemented routers. Future work will select the configuration from
the host's state carrier, test further trigger combinations, and
validate against field or trace-driven current data.
